\section{3. Reihen}

\Def $\sum_{i=0}^{\infty} a_i$ heisst \key{Reihe},
$a_n$ \key{Glieder}/\key{Elemente}/\key{Summanden}

\Def Falls Folge der \key{Teilsummen}/\key{Partialsummen}
$s_n = \sum_{k=0}^{n} a_k$ gegen Grenzwert $L < \infty$ konvergiert,
Reihe \key{konvergiert}, $\lim\limits_{n\to\infty} s_n = \sum_{k=0}^{\infty} a_k = L$,
sonst \key{divergent}. $L$ ist \key{Wert} der Reihe

\Satz \key{Wichtige Reihen}
\vspace{-0.6em}
\begin{itemize}
    \item \key{Geometrische Reihe}: $a_k = a_0 q^k$. Partialsumme $s_n = \sum_{k=0}^{n} a_0 q^k = a_0\frac{1-q^{n+1}}{1-q}$ ($q\ne 1$).
    Für $|q|<1$: $\sum_{k=0}^{\infty} a_0 q^k = \frac{a_0}{1-q}$, für $|q|\ge 1$ divergent
    \item \key{Arithmetische Reihe}: $a_k = a_0 + kd$. Partialsumme
    $s_n = \sum_{k=0}^{n} a_k = (n+1)a_0 + d\frac{n(n+1)}{2} = \frac{(n+1)(a_0+a_n)}{2}$;
    divergiert für $n\to\infty$ (ausser $a_0=d=0$)
    \item \key{Verallg. harmonische Reihe} ($p$-Reihe): $\sum_{n=1}^{\infty} \frac{1}{n^p}$ konvergiert $\iff p > 1$
    (divergiert für $p\le 1$); geschlossene Form i.A. unbekannt, ausser z.B. $p=2$: $\sum \frac{1}{n^2} = \frac{\pi^2}{6}$
    \item \key{Teleskopreihe}: $\sum_{k=1}^{n} (a_k - a_{k+1}) = a_1 - a_{n+1}$.
    Falls $a_n \to L$: $\sum_{k=1}^{\infty} (a_k - a_{k+1}) = a_1 - L$
\end{itemize}

\Satz \key{Notwend. Kriterium für Konvergenz}: \key{Nullfolge}, $\lim\limits_{n\to\infty} a_n = 0$

\Satz $\sum_{n=0}^{\infty} a_n$, $\sum_{n=0}^{\infty} b_n$ konvergent, $C \in \RR$
\vspace{-0.6em}
\begin{itemize}
    \item $\sum_{n=0}^{\infty} (a_n + b_n) = \sum_{n=0}^{\infty} a_n + \sum_{n=0}^{\infty} b_n$
    \item $\sum_{n=0}^{\infty} C \cdot a_n = C \cdot \sum_{n=0}^{\infty} a_n$
\end{itemize}

\Satz Für $N \in \NN_0$, $\sum_{n=N}^{\infty} a_n$ konvergent $\iff$ $\sum_{n=0}^{\infty} a_n$ konvergent.
In diesem Fall gilt $\sum_{n=0}^{\infty} a_n = \sum_{n=0}^{N-1} a_n + \sum_{n=N}^{\infty} a_n$

\Satz Für $a_n\ge 0$ ist $s_n$ wachsend; Reihe konvergiert $\iff$ $s_n$ beschr.

\Satz \key{Vergleichskriterium}: Reihen $a_n, b_n, c_n\ge 0$,
$\exists N \in \NN \, \forall n\ge N : c_n\le b_n\le a_n$. Dann gilt:
\vspace{-0.6em}
\begin{itemize}
    \item $\sum_{n=0}^{\infty} a_n$ konver. $\implies$ $\sum_{n=0}^{\infty} b_n$ konver. (\key{Majorantenkriterium})
    \item $\sum_{n=0}^{\infty} c_n$ diverg. $\implies$ $\sum_{n=0}^{\infty} b_n$ diverg. (\key{Minorantenkriterium})
\end{itemize}

\Def $\sum_{n=0}^{\infty} a_n$ \key{absolut konvergent} falls $\sum_{n=0}^{\infty} |a_n|$ konvergent.
Falls konvergent aber \textbf{nicht} absolut konvergent, \key{Bedingt konvergent}

\Satz \key{Riemannscher Umordnungssatz}: Bedingt konvergent und beliebiges $L \in \RR$:
$\exists \, \varphi: \NN_0 \to \NN_0$ bijektiv mit $\sum_{n=0}^{\infty} a_{\varphi(n)} = L$

\Satz \key{Umordnungssatz für absolut konv. Reihen} Absolut konv., $\varphi: \NN_0 \to \NN_0$ bijektiv.
$\sum_{n=0}^{\infty} a_{\varphi(n)} = \sum_{n=0}^{\infty} a_n$ konvergiert auch absolut

\Satz \key{Leibniz Kriterium}: monoton fallende Nullfolge, $a_n \ge 0$.
Die \key{alternierende Reihe} $\sum_{n=0}^{\infty} (-1)^n a_n$ konvergiert und es gilt:
$\sum_{n=0}^{2m+1} (-1)^n a_n\le \sum_{n=0}^{\infty} (-1)^n a_n\le \sum_{n=0}^{2m} (-1)^n a_n$ für alle $m \in \NN_0$

\Satz \key{Cauchy Kriterium}: Reihe konvergiert $\iff$ für alle $\varepsilon > 0$ existiert ein $N \in \NN_0$ so dass für $n\ge m\ge N$ gilt $|\sum_{k = m+1}^{n} a_k|\le \varepsilon$

\Satz Jede absolut konvergente Reihe konvergiert

\Satz \key{Verallgemeinerte Dreiecksungleichung}: $\left|\sum\limits_{n=0}^{\infty} a_n \right|\le \sum\limits_{n=0}^{\infty}|a_n|$

\Satz \key{Wurzel-Kriterium}: $\rho = \limsup\limits_{n\to\infty} \sqrt[n]{|a_n|} \in \RR \cup \{\infty\}$
\vspace{-0.6em}
\begin{itemize}
    \item $\rho < 1 \implies \sum_{n=0}^{\infty} a_n$ konvergiert absolut
    \item $\rho > 1 \implies \sum_{n=0}^{\infty} a_n$ divergiert
\end{itemize}

\Satz \key{Quotient-Kriterium}: $a_n\ne 0$, $\rho = \lim\limits_{n\to\infty} \frac{|a_{n+1}|}{|a_n|}$
\vspace{-0.6em}
\begin{itemize}
    \item $\rho < 1 \implies \sum_{n=0}^{\infty} a_n$ konvergiert absolut
    \item $\rho > 1 \implies \sum_{n=0}^{\infty} a_n$ divergiert
\end{itemize}

\Satz \key{Cauchy-Produkt}: $a_n, b_n$ absolut konvergent, dann gilt:
$(\sum_{n=0}^{\infty} a_n)(\sum_{n=0}^{\infty} b_n) = \sum_{n=0}^{\infty} (\sum_{k=0}^{n} a_{n-k} b_k)$,
rechts absol. konverg.

\Def Reihe der Form $\sum_{k=0}^{\infty} c_k(x-a)^k$ heisst \key{Potenzreihe},
$a$ heisst \key{Entwicklungspunkt}, $c_0, c_1, \ldots$ heissen \key{Koeffizienten}, $x$ heisst \key{Argument}

\Def Falls, für festes $x$, $s_n(x) = \sum_{k=0}^{n} c_k (x-a)^k$ konvergiert,
Potenzreihe \key{konvergiert} für Argument $x$

\Satz
\vspace{-0.6em}
\begin{itemize}
    \item Jede Potenzreihe hat \key{Konvergenzradius} $R$. Reihe konvergiert für $|x-a| < R$, divergiert für $|x-a| > R$
    \item Intervall $(a-R, a+R)$ heisst \key{Konvergenzintervall}
    \item $\rho = \limsup\limits_{n\to\infty} |c_n|^\frac{1}{n}$. \vspace{-2em}
        \begin{equation*}R = \begin{cases}
            0\text{, falls } \rho = \infty \\
            \rho^{-1}\text{, falls } 0 < \rho < \infty \\
            \infty\text{, falls } \rho = 0
        \end{cases}\end{equation*}
\end{itemize}

\Satz $\star$ $(a_n)_{n\in\NN_0}$ Folge komplexer Zahlen und
$$R = \begin{cases} \frac{1}{\limsup \limits_{n\to\infty}|a_n|^{\frac{1}{n}}} \hspace{10pt}\text{falls } \limsup\limits_{n\to\infty} |a_n|^{\frac{1}{n}} > 0 \\ +\infty \hspace{43pt}\text{sonst}\end{cases}$$
Dann konv. $\sum^\infty_{n=0} a_nz^n$ absolut für $|z| < R$ und diverg. für $|z| > R$

\Satz $\star$ $(a_n)_{n\in\NN_0}$ Folge komplexer Zahlen bei der 0 nicht unendlich oft auftritt und
$$R = \lim_{n\to\infty} \left|\frac{a_n}{a_{n+1}}\right|$$
Falls $R$ existiert konv. $\sum^\infty_{n=0} a_nz^n$ absolut für $|z| < R$ und diverg. für $|z| > R$

\Satz $\star$ $\sum_{n=1}^\infty \frac{1}{n^2} = \frac{\pi^2}{6}$
